\documentclass[11pt,letterpaper]{report}
\input{preambulo}
\begin{document}
\hypersetup{pdftitle={Observaciones al borrador y auditoría R01 - Minicentral Limbani R02}}
\chapter*{Informe de observaciones al borrador}
\textbf{Proyecto:} Minicentral hidroeléctrica Limbani, primera unidad.\\
\textbf{Alumno:} Aquiles Taylor Ramos Yapo, 227368.\\
\textbf{Docente:} MEZA MAROCHO GREGORIO. \textbf{Sección:} GRUPO A.\\
\textbf{Base:} PDF recibido el 13 de octubre de 2026, 28 páginas físicas (27 numeradas y portada).\\
\textbf{Revisión:} R02, técnica, normativa y documental.

Se identifican 24 observaciones. ``Crítica'' significa que cambia la factibilidad, la seguridad o la base de diseño; ``Alta'' requiere resolución antes de especificar o construir; ``Media'' afecta presentación y trazabilidad. Corregir una afirmación en el informe no acredita que el equipo o instalación física cumpla.

Las páginas siguientes se refieren a la numeración impresa del borrador, no al índice físico del lector PDF. Las fuentes N01-N20, las ecuaciones y el registro de evidencias están desarrollados en el expediente R02 y en \texttt{datos/normativa.json}.
\setcounter{section}{0}
\renewcommand{\thesection}{\arabic{section}}
\begingroup
\small
\setlength{\parskip}{3pt}
\titlespacing*{\section}{0pt}{.6em}{.2em}
\input{generado/observaciones}
\endgroup
\clearpage
\chapter*{Auditoría de R01 y correcciones R02}
Además de los 24 hallazgos del borrador original, se identifican nueve mejoras de la revisión anterior. Los resultados antiguos se conservaron por separado y no se usan como salidas activas.
\setcounter{section}{0}
\begingroup\small
\setlength{\parskip}{3pt}
\titlespacing*{\section}{0pt}{.6em}{.2em}
\input{generado/revision-r02}
\endgroup
\end{document}
